\documentclass[spanish,10pt,a5paper]{article}
\usepackage[T1]{fontenc}
\usepackage{graphicx}
\usepackage{xcolor}
\usepackage{babel}
\usepackage{tikz}
\usetikzlibrary{babel,calc}
\usepackage{kpfonts}



% Definición de los colores, es usan colores que queda con la imagene de fondo
\definecolor{col1}{HTML}{C66545}
\definecolor{col2}{HTML}{AC7D73}
\definecolor{col3}{HTML}{7A9584}

\begin{document}
\thispagestyle{empty}

\begin{tikzpicture}[remember picture, overlay]

\node[inner sep=0pt] at (current page.center) {\includegraphics[height=\paperheight]{1 (1).png}};


\node[ inner sep=0pt] at (current page.center) {%
\begin{minipage}[c][0.55\paperheight]{0.65\paperwidth} 

\centering
{\color{col1}\Huge\bfseries Titulo del tema }\\[3mm]
{\color{col2}\huge\itshape Descripción del tema}\\[5mm] 

{\color{col3}\large Nombre del autor para la portada en este espacio }\\[15mm]
\begin{minipage}{0.5\paperwidth} 
{\color{col1!45!black}
\centering {una breve introduccion al tema.}
} 
\end{minipage}\\[11mm]
{\color{col1}\large aqui va el numero de la actividad 
}
\end{minipage}};


\fill[col3] ($(current page.north east)+(-4,4)$) -| ++(4,-4) --cycle;
\node[white, rotate=-45] at ($(current page.north east)+(-1.5,-1.5)$) 
{\large\bfseries fecha}; 
\end{tikzpicture}
\end{document}
